\documentclass[11pt,a4paper]{article}
\usepackage{amsmath,amssymb,amsthm}
\usepackage[margin=2.5cm]{geometry}
\usepackage{hyperref}

\newtheorem{definition}{Definition}[section]
\newtheorem{theorem}[definition]{Theorem}
\newtheorem{proposition}[definition]{Proposition}
\newtheorem{lemma}[definition]{Lemma}
\newtheorem{corollary}[definition]{Corollary}
\newtheorem{conjecture}[definition]{Conjecture}
\newtheorem{remark}[definition]{Remark}

\title{Hyperseed Formalization:\\
  Google's Quantum Crypto Paper Tells Us More Than They Intended}
\author{ProtomegaTron (automated formalization)}
\date{2026-09-18}

\begin{document}
\maketitle

\begin{abstract}
We formalize the intellectual structure of Ben Goertzel's Substack article
``Google's Quantum Crypto Paper Tells Us More Than They Intended''
(2026-04-01) within the Hyperseed ontology. The article analyzes
Google Quantum AI's whitepaper on breaking ECC with Shor's algorithm.
Key mappings: the responsible disclosure paradox as section-fiber leakage,
threat-capability coupling as dual fiber, search space bounding as
fiber constraint propagation, post-quantum migration as fiber transition,
and harvest-now attacks as fiber time asymmetry.
\end{abstract}

\tableofcontents

%% =================================================================
\section{Section-fiber leakage in responsible disclosure}
\label{sec:leakage}
%% =================================================================

\begin{theorem}[The disclosure paradox --- claims 1--4, 18--20, 21]
\label{thm:leakage}
Google withholds the specific circuit implementation (the ``crown jewel'')
while publishing the complete architectural framework: kickmix architecture,
windowed arithmetic, Shor phase estimation with qubit-recycled QFT,
Montgomery's trick, yoked surface codes, reaction-limited execution,
and exact resource counts from ZK proof statements.

In Hyperseed terms, this is \emph{section-fiber leakage}: the withheld
circuit (section) leaks through the published architectural details
(fiber). The fiber constrains the section so tightly that withholding
the section is insufficient:
\[
  \text{Published fiber} \implies \text{bounded search over sections}
  \implies \text{section reconstructible}
\]

The mechanism: the paper specifies enough fiber constraints (gate counts,
qubit counts, architecture type, optimization strategy, design space
boundaries) that the unknown section (specific circuit wiring) is
determined up to a bounded search. Information-theoretically:
\[
  H(\text{circuit} \mid \text{published constraints})
  \ll H(\text{circuit})
\]

The residual entropy --- what you still need to figure out after reading
the paper --- is small enough that a competent team of 3--5 researchers
can search the remaining space in months. This is security through
obscurity, and it fails for the same reason it always fails: the
``secret'' lives in a space that's bounded by publicly known constraints.

This is the same pattern as the Leaky Transcension Hypothesis: perfect
containment is impossible because the fiber projects onto the base.
Google's paper is a transcended interior (the secret circuit) that leaks
unavoidably through its boundary (the published constraints). The leakage
isn't accidental --- it's structural. You can't describe an architecture
precisely enough to be scientifically credible without constraining its
implementations enough to be reproducible.
\end{theorem}

%% =================================================================
\section{Fiber constraint propagation}
\label{sec:propagation}
%% =================================================================

\begin{proposition}[Bounded reverse engineering --- claims 5--8, 23]
\label{prop:propagation}
The paper's disclosed constraints propagate to bound the unknown
components. The design space for elliptic curve point addition on
a surface-code quantum computer is well-studied: modular arithmetic
strategy, coordinate system selection, and ancilla management are
textbook topics. Prior work (Litinski 2023, H\"aner et al.\ 2020,
Gouzien et al.\ 2023, Chevignard et al.\ 2026) establishes the
space; Google improved spacetime volume ${\sim}10\times$ within it.

In Hyperseed terms, this is \emph{fiber constraint propagation}:
known fiber constraints on a system determine its unknown components
up to a bounded search:
\begin{enumerate}
\item \textbf{Architecture type} (kickmix, windowed) constrains the
  class of circuits.
\item \textbf{Resource counts} (gate counts, qubit counts) constrain
  the scale.
\item \textbf{Prior work} constrains the design space (known
  optimizations, known tradeoffs).
\item \textbf{Performance claims} constrain the quality (must achieve
  the stated spacetime volume).
\end{enumerate}

Each constraint eliminates a region of the search space. Together,
they leave a bounded residual --- the specific circuit wiring ---
that is a constrained optimization problem, not an open research
question.

This is a general principle: enough fiber constraints on any system
determine its unknown components. The more fiber you publish, the
less section remains hidden. Google published enough fiber to make
the section reproducible.
\end{proposition}

%% =================================================================
\section{Dual fiber: threat and capability}
\label{sec:dual}
%% =================================================================

\begin{theorem}[Threat-capability coupling --- claims 12--14, 22]
\label{thm:dual}
The qubit counts for cryptographic attacks (${\sim}500$K physical)
are roughly the same as for quantum-enhanced AI. The threat and the
capability arrive on the same timescale, on the same hardware, driven
by the same engineering progress.

In Hyperseed terms, this is \emph{dual fiber}: the same physical fiber
(qubit count, error correction architecture, gate fidelity) supports
both threat (breaking cryptography) and capability (quantum-enhanced
AI). They are inseparable aspects of the same underlying fiber:
\[
  F_{\text{quantum}} = F_{\text{threat}} \cup F_{\text{capability}}
  \qquad \text{with} \qquad
  F_{\text{threat}} \cap F_{\text{capability}} \neq \emptyset
\]

The intersection is substantial: the same logical qubits, the same
error correction, the same gate operations. You can't build quantum
hardware capable of one without being capable of the other. The dual
use is not a policy choice but a physical constraint: the fiber is
inherently dual.

This means that any attempt to prevent the threat (restrict quantum
cryptanalysis) necessarily restricts the capability (quantum AI).
And any progress on the capability side necessarily advances the
threat side. The policy implications are that threat-only analysis
misses half the picture --- arguably the more important half.
\end{theorem}

%% =================================================================
\section{Fiber transition: post-quantum migration}
\label{sec:transition}
%% =================================================================

\begin{definition}[Cryptographic fiber transition --- claims 15--17, 24--25]
\label{def:transition}
Every blockchain using standard ECC is on a countdown. Migration to
post-quantum cryptography is urgent --- not a future concern but a
current engineering priority.

In Hyperseed terms, this is a \emph{fiber transition}: replacing the
cryptographic fiber of the entire blockchain ecosystem:
\[
  F_{\text{ECC}} \xrightarrow{\text{migration}} F_{\text{PQ}}
  \qquad \text{before} \qquad
  F_{\text{quantum}} \text{ matures}
\]

It's a race between fiber transitions: the defensive transition
(ECC to post-quantum) must complete before the offensive transition
(quantum computers reaching cryptanalytic capability) arrives.

\textbf{Harvest-now attacks} introduce a \emph{fiber time asymmetry}:
data encrypted under the ECC fiber today becomes vulnerable when the
quantum fiber arrives later. The commitment (encryption) is irreversible
--- the ciphertext is already public on the blockchain. But the security
guarantee expires when the decryption fiber (quantum computation) becomes
available:
\[
  \text{Security}(t) = \begin{cases}
    \text{high} & \text{if } F_{\text{quantum}}(t) < \text{threshold} \\
    \text{zero} & \text{if } F_{\text{quantum}}(t) \geq \text{threshold}
  \end{cases}
\]

The discontinuity is what makes it urgent: security doesn't degrade
gradually --- it collapses to zero once the quantum fiber crosses the
threshold. And the threshold crossing is approaching faster than most
blockchain projects are migrating.
\end{definition}

%% =================================================================
\section{Cross-references}
%% =================================================================

\begin{remark}[Connections to other formalizations]
\begin{itemize}
\item \textbf{The Leaky Transcension Hypothesis} (2026-04-25):
  section-fiber leakage. The same pattern --- perfect containment
  is impossible because fiber constrains section --- applied to
  information disclosure rather than civilizational transcension.
\item \textbf{DeAI After Mythos} (2026-04-15): security regime change.
  Mythos changes the economics of vulnerability discovery; quantum
  changes the economics of cryptographic breaking. Both are
  base-level security collapses.
\item \textbf{The OpenAI/Microsoft Shakeup} (2026-05-01): structural
  constraints. The disclosure paradox parallels the governance
  paradox: you can't structurally depend on secrecy any more than
  you can structurally depend on good intentions.
\item \textbf{Toward a Truly Decentralized Digital} (2026-06-25):
  blockchain security. Post-quantum migration is essential for
  the decentralized digital identity infrastructure.
\item \textbf{Bootleggers \& Baptists} (2026-04-28): dual use
  regulation. The dual-fiber nature of quantum computing means
  that regulation targeting the threat side necessarily affects
  the capability side.
\end{itemize}
\end{remark}

\begin{conjecture}[Fiber disclosure threshold]
Conjecture: for any technical system, there exists a disclosure threshold
--- a minimum amount of fiber that must be published for scientific
credibility --- above which the remaining section is reconstructible
by competent researchers. Google's paper crosses this threshold.

If this conjecture holds, ``responsible disclosure'' of technical
achievements faces a fundamental tension: scientific credibility
requires publishing enough fiber to be verifiable, but verifiable
fiber constrains the section enough to be reproducible. The only
way to truly withhold a capability is to not publish at all ---
and even then, independent discovery is possible.

This suggests that for dual-use quantum advances, the choice is not
between disclosure and secrecy but between acknowledged disclosure
(with coordinated response) and delayed independent discovery (without
coordination). Acknowledged disclosure is generally safer.
\end{conjecture}

\end{document}
